\chapter{Introduction}
\label{chap:intro}
\minitoc

\section{Exemple d'illustration}

\subsection{Une sous section pour le fun}

Eh non je n'�crirai pas votre th�se... D�sol�. Juste pour montrer que ce style g�re les accents et comment on inclut une figure. A noter, les path vers les dossiers d'images et l'extension sont d�finis dans formatAndDefs.tex . Pour pdflatex, les images jpg sont g�r�es. Pour latex, utiliser ps ou eps.

\begin{figure}[!htbp]
  \begin{center}
    \includegraphics[width=0.9\textwidth]{Chapitre1/arctic_control}
  \end{center}
  \caption{Une jolie image...}
  \label{fig:jolieImage}
\end{figure}

\section{Une �quation}

Juste pour montrer les commandes argmin et d�riv�e partielle.

\begin{equation}
  T = \argmin_T E(T,R,F)
\end{equation}

R�gularisation:

\begin{equation}
  \pd{T}{t} = \Delta T
\end{equation}

\section{Une autre section}

Pour montrer l'environnement bulletList

\begin{bulletList}
 \item Premier point
 \item Deuxi�me point
% \item Et une r�f�rence ajout�e dans la liste des abr�viations \nomenclature{DTI}{Diffusion Tensor Imaging} DTI
\end{bulletList}
